\begin{figure}[htbp]
\centering
\begin{tikzpicture}[font=\small,>=Stealth]
\fill[blue!10,even odd rule] (0,0) circle(2) (0,0) circle(.95);
\draw[thick] (0,0) circle(2);\draw[thick] (0,0) circle(.95);
\draw[blue!70!black,very thick,->] (20:2) arc(20:95:2);
\draw[blue!70!black,very thick,->] (200:.95) arc(200:120:.95);
\draw[red!75!black,thick,->] (2,0)--(2.8,0) node[right] {outward};
\draw[red!75!black,thick,->] (.95,0)--(.25,0);
\node[font=\footnotesize] at (0,-.18) {hole};
\node[blue!70!black,above] at (0,2.2) {outer: counterclockwise};
\node[blue!70!black,below] at (0,-2.2) {inner: clockwise};
\draw[->] (-3.1,-.3)--(-2.45,-.3) node[right] {$e_1$};
\draw[->] (-3.1,-.3)--(-3.1,.4) node[above] {$e_2$};
\node[align=left] at (4,1.25) {surface orientation:\\$\dd x\wedge\dd y$};
\node[align=left] at (4,-1.1) {$(\nu,\tau)$ positive:\\outward vector first,\\boundary tangent second};
\end{tikzpicture}
\caption{Outward-first on a standard oriented annulus. At the rightmost
outer point, outward is right and the positive tangent is up. At the
rightmost inner point, outward is left into the hole and the positive
tangent is down. Thus the two boundary components run in opposite senses.}
\label{fig:annulus-orientation}
\end{figure}
